\section{Rigorous Foundations: Three Central Results}
\label{sec:rigorous-vector-foundations}

The geometric rules developed in this chapter rest on a small number of precise
inequalities and decomposition theorems. Establishing them now clarifies why
projections are well defined and why Euclidean length behaves as expected.

\subsection{Cauchy--Schwarz Inequality}
For vectors $\mathbf u$ and $\mathbf v$ in a real inner-product space,
\[
\boxed{|\mathbf u\cdot\mathbf v|\leq \lVert\mathbf u\rVert\,\lVert\mathbf v\rVert.}
\]
If $\mathbf v=\mathbf0$ the statement is immediate. Otherwise consider, for
real $t$, the nonnegative squared norm
\[
0\leq\lVert\mathbf u-t\mathbf v\rVert^2
=\lVert\mathbf u\rVert^2-2t(\mathbf u\cdot\mathbf v)
+t^2\lVert\mathbf v\rVert^2.
\]
Choose $t=(\mathbf u\cdot\mathbf v)/\lVert\mathbf v\rVert^2$. Substitution gives
\[
0\leq \lVert\mathbf u\rVert^2-
\frac{(\mathbf u\cdot\mathbf v)^2}{\lVert\mathbf v\rVert^2},
\]
and multiplication by $\lVert\mathbf v\rVert^2$ yields the result. Equality
holds precisely when $\mathbf u=t\mathbf v$, so the vectors are linearly
dependent only when the inequality is strict.

\begin{physicalbox}[Geometric intuition]
Cauchy--Schwarz states that the signed shadow of one vector on another can
never be longer in magnitude than the original vector. It is the rigorous
reason that $|\cos\theta|\leq1$ in the inner-product formula.
\end{physicalbox}

\begin{mistakebox}[Common misconception: the inequality is usually strict]
Equality is not evidence that two vectors are equal. It means that they are
linearly dependent: they may point in the same direction or in opposite
directions, and they may have different lengths.
\end{mistakebox}

\subsection{Triangle Inequality}
The length of a sum cannot exceed the sum of the lengths:
\[
\boxed{\lVert\mathbf u+\mathbf v\rVert\leq
\lVert\mathbf u\rVert+\lVert\mathbf v\rVert.}
\]
Indeed,
\begin{align*}
\lVert\mathbf u+\mathbf v\rVert^2
&=\lVert\mathbf u\rVert^2+2\mathbf u\cdot\mathbf v+\lVert\mathbf v\rVert^2\\
&\leq\lVert\mathbf u\rVert^2+2\lVert\mathbf u\rVert\lVert\mathbf v\rVert+\lVert\mathbf v\rVert^2\\
&=(\lVert\mathbf u\rVert+\lVert\mathbf v\rVert)^2,
\end{align*}
where Cauchy--Schwarz was used in the second line. Taking nonnegative square
roots proves the claim.

\begin{physicalbox}[Geometric intuition]
Travelling directly from the initial point to the final point is never longer
than travelling through an intermediate point. Equality occurs when the two
displacements are aligned in the same direction.
\end{physicalbox}

\subsection{Orthogonal Decomposition Theorem}
Let $\mathbf b\neq\mathbf0$. Every vector $\mathbf a$ has a unique decomposition
\[
\boxed{\mathbf a=\mathbf a_{\parallel}+\mathbf a_{\perp}},
\qquad
\mathbf a_{\parallel}\parallel\mathbf b,
\quad
\mathbf a_{\perp}\perp\mathbf b.
\]
Existence follows by defining
\[
\mathbf a_{\parallel}=\frac{\mathbf a\cdot\mathbf b}{\mathbf b\cdot\mathbf b}\mathbf b,
\qquad
\mathbf a_{\perp}=\mathbf a-\mathbf a_{\parallel}.
\]
The dot product $\mathbf a_{\perp}\cdot\mathbf b$ vanishes directly. For
uniqueness, suppose
$\mathbf a=\mathbf p_1+\mathbf q_1=\mathbf p_2+\mathbf q_2$, where each
$\mathbf p_i$ is parallel to $\mathbf b$ and each $\mathbf q_i$ is perpendicular
to it. Then
\[
\mathbf p_1-\mathbf p_2=-(\mathbf q_1-\mathbf q_2).
\]
The left side is parallel to $\mathbf b$ and the right side perpendicular to
$\mathbf b$. The only vector with both properties is $\mathbf0$, hence
$\mathbf p_1=\mathbf p_2$ and $\mathbf q_1=\mathbf q_2$.

\begin{mistakebox}[Common misconception: projection requires a unit vector]
The formula $\operatorname{proj}_{\mathbf b}\mathbf a=(\mathbf a\cdot\mathbf b)\mathbf b$
is valid only when $\lVert\mathbf b\rVert=1$. For a general nonzero $\mathbf b$,
the denominator $\mathbf b\cdot\mathbf b$ is essential.
\end{mistakebox}

\subsection{Why These Results Matter Later}
Cauchy--Schwarz controls probabilities and expectation values in quantum
mechanics. The triangle inequality underlies every metric and error estimate.
Orthogonal decomposition becomes Fourier expansion, normal-mode analysis,
least-squares approximation, and projection onto quantum measurement outcomes.
